\section{Introduction and Motivation}
The main goal of representation theory is to understand abstract algebraic structures—specifically groups and associative algebras—by concretely realizing their elements as linear transformations of vector spaces. By transforming non-linear group multiplication into the well-understood domain of linear algebra and matrix arithmetic, one converts difficult group-theoretic questions into tractable eigenvalue problems, trace computations, and spectral decompositions.
\\
First we define what is a representation? then we move to defining a homomorphism and isomorphism of representations. Moreover, we introduce the definition of an important structure that we will be using extensivly throughout this exposition which is the group algebra $k[G]$, then we will move to working with simple representations/modules. In between we show how can modules and representations be seen as the same thing.

\section{Basic Definitions}

\begin{definition}[Representation of an Algebra]
    
    Let $A$ be a $k-$algebra, and $V$ be a $k-$vector space. A \emph{representation} of $A$ is a pair $(V,\rho)$ where $\rho: A \to \operatorname{End}(V)$ is a $k-$algebra homomorphism.
\end{definition}
\begin{definition}[Subrepresentation]
    
    Let $(V, \rho)$ be a representation of any $k-$algebra $A$. We say that a subspace $W \subseteq V$ is a subrepresentation iff elements of $W$ commute with the representation $rho$ or that $W$ is invariant under the action of $A$. That is for any $a \in A$ and any $w \in W$ we have that 
    \[
        \rho(a) \cdot w \in W
    \]
    
\end{definition}

\begin{definition}[Intertwining Operator]
    Let $(V, \rho_1)$ and $(W, \rho_2)$ be any two representations of a $k-$algebra $A$, an intertwining operator $T:V \to W$ is a linear map that commutes with the action of $A$, that is
    \[
       T(\rho_1(a)(v)) = \rho_2(a)(T(v))
    \]
    We may also say that $T$ is a homomorphism of representations. If $T$ is invertible, then it is an isomorphism of representations
\end{definition}


\subsection{Equivalence of Representations and Modules}
Now, how are representations of any $k-$algebra $A$ could be seen as $A$-modules?

We basically define the action of $A$ on $V$ by the action of the image of $A$ in $\text{End}(V)$ as linear transformations.
 For the other direction, Let $M$ be an $A-$module. Then $M$ is an abelian group. We define scalar multiplication via the action of $A$, as we mentioned earlier $A$ has a copy of the field $k$ inside it, hence: for any $\lambda \in k$ and any $m \in M$ we have that
 \[
 \lambda \cdot  m = \lambda m
 \]w
 where $\cdot$ is the action of $A$ on $M$. It remains to prove that the action of $A$ on $M$ as a module defines group algebra homomorphism from $A$ to $\text{End}(M)$.

\begin{proof}
    Fix any $a \in A$, consider the map
    \[
    \varphi(a): M \to M
    \]
    given by $m \mapsto a \cdot m$. By the axioms of a module action we get 
    \begin{itemize}
        \item $\varphi(a) (m+n) = a \cdot (m + n) = a \cdot m + a \cdot n $
        \item $\varphi(a) (\lambda m) = a \cdot (\lambda m) = (a \lambda) \cdot m = (\lambda a) \cdot (m) = \lambda(\varphi(a) (m)) 
         $
    \end{itemize}
    the last equality is because we defined the image of the field $k$ in $A$ to be central. So, we proved that $\varphi(a)$ is a linear map for any $a \in A$. Now, we define the group algebra homomorphism
    \[
    \rho: A \to \text{End}(M)
    \]
    given by $a \mapsto \varphi(a)$. To see that it is an algebra homomorphism, first we use the linearity of $\varphi(a)$
    \begin{itemize}
    \item $\rho(a + b ) = \varphi(a+b) = \varphi(a) + \varphi(b) = \rho(a) + \rho(b)$
    \item $\rho( \lambda a) = \varphi(\lambda a) = \lambda \varphi(a) = \lambda \rho(a)$
    
    \end{itemize}
    since $M$ is an $A$ module, then by associativity of the action and the definition of $\varphi(a)$ for any $a \in A$, we get that 
\begin{itemize}
    \item $\rho(ab) = \varphi(ab) = \varphi(a) \varphi(b) = \rho(a)\rho(b)$
\end{itemize}

\end{proof}
\section{The Group Algebra}
Now, we move to a special type of algebra called the Group algebra $k[G]$, that will help us study the representations of a finite group $G$.


\begin{definition}[The Group Algebra $k[G]$]

Let $G$ be a finite group and $k$ a field. The \emph{group algebra} $k[G]$ is the $k$-vector space with basis given by the elements of $G$:

\[
k[G] = \left\{ \sum_{g \in G} a_g \, g \;\middle|\; a_g \in k \right\}.
\]

The multiplication on $k[G]$ is defined by extending the group multiplication linearly:

\[
\left( \sum_{g \in G} a_g \, g \right) \left( \sum_{h \in G} b_h \, h \right) = \sum_{g, h \in G} a_g b_h \, (gh) = \sum_{x \in G} \left( \sum_{gh = x} a_g b_h \right) x.
\]

The vector space dimension of $k[G]$ over $k$ is precisely the order of the group:

\[
\dim_k(k[G]) = |G|.
\]

\end{definition}

Before we define a representation of a group, we define the general case of a representation of a group algebra.

\begin{definition}[Representation of a group algebra $k[G]$]
    Let $G$ be a group. Let $k[G]$ be any group algebra over a field $k$. A representation of $k[G]$ is a pair $(V, \rho)$, where $V$ is a $k$-vector space and
   \[
        \rho: A \to \text{End}(V)
    \]
    
    is a group algebra homomorphism.

\end{definition}

The previous definition just says that we are trying to represent elements of an algebra with linear transformation or matrices.






\section{Representation of a Group}
We are done discussing the general case. But what about the definition of a group representation? 
There are two ways to view this definition. One way is to see it as a restriction of the definition of an algebra representation on the basis elements which are basically the elements of the group $G$. The other one, which I find more interesting is to see it as an extension of a group action.

First, recall the definition of a group action on an arbitrary set $S$.

\begin{definition}[Group Action]
    Let $G$ be a multiplicative group, and $S$ to be any set. We define a group action of $G$ on $S$ to be a function 
    \[
    \cdot: G \times S \to S
    \]
    denoted by $g \cdot x$ for all $g \in G$ and $x \in S$, satisfying the following axioms:
    \begin{enumerate}
        \item $e \cdot x = x  \quad\forall \,\, x \in S$
        \item $g \cdot (h \cdot x) = gh \cdot x$
    \end{enumerate}

\end{definition}
Using the previous two axioms of a group action, one can see that it actually induces a group homomorphism $\alpha: G \to \operatorname{Sym}(S)$, where $\operatorname{Sym}(S)$ is the group of all permutations on $S$ with function composition.
\begin{proof}
    Let $\cdot : G \times S \to S$ be a group action. Consider the map 
    \[
    \alpha: G \to \operatorname{Sym}(G)
    \] 
    where $\alpha(g)$ is a bijection on $S$ such that 
    
    \[\alpha(g)(x) = g \cdot x \]
    for all $g \in G$ and $x \in S$. We claim that $\alpha$ is a group homomorphism. By the axioms of the group action, for any $g,h \in G$: 
    \[
    \alpha(gh)(x) = gh \cdot x = g \cdot (h \cdot x) = \alpha(g)(\alpha(h)(x)) = \alpha(g) \circ \alpha(h) (x)
    \]
\end{proof}

Now, we repeat the same process but using a group action of a group $G$ on a vector space $V$ instead of an arbitrary set $S$ requiring an additional property for the group action which is to be $\emph{linear}$.

\begin{definition}[Linear Group Action]
    Let $V$ be a $k-$vector space, and $G$ be any group. Define the group action of $G$ on $V$ as the function 
    \[
    \cdot: G \times V \to V
    \]
    that satisfies the following axioms:
    \begin{enumerate}
        \item $e \cdot v = v  \quad\forall \,\, v \in V$
        \item $g \cdot (h \cdot v) = gh \cdot v$
        \item $g \cdot (v + w) = g \cdot v + g \cdot w$
        \item $g \cdot ( \lambda v) = \lambda (g \cdot v)$
    \end{enumerate}
\end{definition}
This definition is in fact equivalent to having a group homomorphism $\varphi: G \to \operatorname{GL}(V)$, exactly in the same way a regular group action gives rise to a homomorphism of groups from $G$ to the symmteric group on an arbitrary set $S$. This pair of the homomorphism $\varphi$ along with the vector space $V$ are called a $\emph{representation}$ of the group $G$.
\begin{definition}[Representation of a Group]

Let $G$ be a group, and let $V$ be a vector space over a field $k$. A \emph{representation} of $G$ is a pair $(V, \varphi)$ where $\varphi: G \to \operatorname{GL}(V)$ is a group homomorphism given by the linear action of $G$ on $V$.
\end{definition}

\begin{definition}[$G$-Homomorphism]

Let $(\rho_1, V_1)$ and $(\rho_2, V_2)$ be two representations of a group $G$ over $k$. A linear map $T \colon V_1 \to V_2$ is called an \emph{intertwining operator} (or \emph{$G$-homomorphism}) if it commutes with the $G$-action:

\[
T \circ \rho_1(g) = \rho_2(g) \circ T \quad \text{for all } g \in G.
\]

 If $T$ is an invertible linear isomorphism, we say that $\rho_1$ and $\rho_2$ are \emph{equivalent} (or \emph{isomorphic}), denoted $\rho_1 \cong \rho_2$.
 % In matrix terms, fixing bases for $V_1$ and $V_2$, there exists an invertible matrix $C$ such that

% \[
% \rho_2(g) = C \, \rho_1(g) \, C^{-1} \quad \text{for all } g \in G.
% \]

\end{definition}

\begin{definition}[Subrepresentation]

Let $\rho \colon G \to \operatorname{GL}(V)$ be a representation. A vector subspace $W \subseteq V$ is called a \emph{$G$-invariant subspace} (or \emph{subrepresentation}) if

\[
\rho(g)w \in W \quad \text{for all } g \in G \text{ and all } w \in W.
\]

\end{definition}

\begin{definition}[Irreducible representation]
    We say a representation $(V, \rho)$ is irreducible if $V$ has no proper non-zero subrepresentations.
\end{definition}
\begin{definition}[Indecomposable representation]
    We say a representation $(V, \rho)$ is indecomposable if there doesn't exist a pair of subrepresentations $U,W$ such that $V = U \oplus W$.
\end{definition}
\begin{note}

It is essential to distinguish between \emph{irreducibility} and \emph{indecomposability}:

\begin{itemize}

    \item A representation $V$ is \emph{irreducible} if it admits no proper non-zero subrepresentation.

    \item A representation $V$ is \emph{indecomposable} if it cannot be written as a direct sum of two non-zero subrepresentations, i.e., $V \neq W_1 \oplus W_2$ with $W_1, W_2 \neq \{0\}$.

\end{itemize}

Every irreducible representation is indecomposable, but the converse does not hold in general fields! As we shall see in Section~\ref{sec:modular_failure}, when the characteristic of the field divides the group order, there exist indecomposable representations that are strictly reducible.

\end{note}

% When analyzing representations, the first foundational question that emerges is: \emph{When should two representations be considered the same?} As observed in elementary linear algebra, a single linear transformation possesses infinitely many matrix representations depending on the chosen basis. Two matrix representations are equivalent if they differ merely by an invertible change-of-basis matrix.

